% chapters/05_discussion.tex - 第5章 讨论
% 结构：按研究目的1→2→3→4对应
% 5.2 对应目的1；5.3 对应目的2
% 5.4-5.6 对应目的3；5.7 对应目的4
% 5.8 局限性；5.9 后续RCT建议

\chapter{讨论}

\section{主要发现概述}
\label{sec:ch5_intro}

本研究围绕自研移动端计算机视觉系统SportSci Pro的测量效度验证及其在数智化监控下肢抗阻训练中的应用, 采用解释性序贯混合方法设计, 系统收集了量化实验数据与质性访谈资料。综合第4章的全部实证结果, 本研究提炼出如下七项核心发现, 并在此基础上明确本预试验性研究的证据层级与科学定位。

目的1（工具效度）的核心发现：第一, 测量学基础稳固。研究一在受控实验条件下验证了SportSci Pro与GymAware在杠铃后蹲平均向心速度测量上具有优秀的一致性（ICC=0.940, MAE=0.047 m/s）；研究二在真实训练场景中的Rep级配对进一步获得ICC=0.991（Bias=+0.0019 m/s）与热身级ICC=0.987的高度一致性证据。然而, 热身阶段存在+0.0212 m/s的正向偏移, 且88对配对100\%均为正偏。

目的2（干预可行性）的核心发现：第二, 训练方案在依从受试者中具备较高可行性, 但招募保留与系统可用性存在改进空间。PP样本24名受试者的课次层训练完成率达98.4\%（189/192）, 高依从率达100\%（24/24）, AI辅助组速度目标命中率达80.2\%$\pm$6.9\%, 自动减组触发忠实度达100\%, 严重不良事件为零。然而, 随机化后主要后测完成率仅为66.7\%（24/36）, AI辅助组SUS评分65.00$\pm$4.61分, 两项预设推进标准未达标。

目的3（干预效果）的核心发现：第三, 两组外部训练剂量高度对等, 但负荷组织形态呈现与反馈路径逻辑相一致的结构分化。AI辅助组总重复次数显著更少（100.3 vs. 105.5次, $p$=0.035, $g$=$-$0.95）, 派生的每次重复平均机械负荷高出约6.1\%。第四, Hooper主观恢复状态在两组间高度平行（组别$p$=0.421, 交互$p$=0.945）, 但课次主效应显著（$p$=0.004）, 两组均表现出先升后降的时间演变形态。第五, 探索性体能适应未表现确证性组间优效, 但CMJ捕获到值得关注的方向性信号（调整后差值+2.16 cm, $p$=0.051, $g$=0.85, $g$ 95\% CI：0.02至1.69）。第六, 训练自我效能呈现极显著的组间优势（调整后差值+9.88分, $p<0.001$, $g$=2.70）。

目的4（用户接受度）的核心发现：第七, 系统可用性评分未达预设推进阈值, 但技术接受度极高, 呈现原型系统的典型二元特征。AI辅助组SUS得分均值65.00$\pm$4.61分（仅P004一人达70分）；感知有用性、信任度与使用意愿三项TAM维度均分均超过4.4分。质性访谈归纳出5类抗阻训练行为特征模式, 揭示了用户自主权与算法控制之间的动态博弈。

上述七大发现共同指向一个核心结论：当前4周预试验在完成流程可行性验证的同时, 尚未能够提供数智化监控相较于规范化RIR自我指导在生理适应层面的确证性优效证据, 但已在心理适应层面捕获到强烈且稳健的优势信号, 并在训练执行行为学层面确认了与速度损失阈值逻辑相一致的结构分化。鉴于保留率（66.7\%）与系统可用性（SUS 65.00分）未达预设推进标准, 当前方案在进入大规模确证性随机对照试验前, 需要实施针对性的招募流程优化与系统交互重构。

\section{移动端视觉监控的测量学效度与技术应用边界}
\label{sec:validity_discussion}

本节回应目的1, 对第4章4.1与4.2呈现的双阶段效度证据进行机制解读与技术应用边界的划定。

\subsection{从实验室效度到现场生态效度的成功迁移}
\label{sec:lab_to_field}

本研究通过双阶段效度验证设计, 系统检验了SportSci Pro移动端视觉监控系统从受控实验环境到真实训练场景的测量学迁移性。研究一在实验室条件下报告了SportSci Pro与GymAware在杠铃后蹲平均向心速度测量上的一致性达到ICC=0.940（MAE=0.047 m/s）, 构成了系统的基线测量学证据；研究二则通过静默同步记录方案, 进一步获取了真实训练场景下的生态效度数据。

Rep级配对分析（n=43）显示的效度水平不仅未因场景迁移而衰减, 反而呈现出优于实验室条件的一致性表现：ICC从0.940上升至0.991, Bias由+0.002 m/s改善至+0.0019 m/s, MAE从0.047 m/s减小至0.017 m/s（缩减约64\%）。

从测量原理层面分析, Rep级现场一致性的显著提升可能源于三个协同因素。首先, 研究二深蹲主项工作组的典型速度区间（0.4--0.8 m/s）恰好位于MediaPipe姿态估计算法与GymAware线性位置传感器的共同最优测量区间。其次, 研究二的训练操作由经过预实验校准的研究员现场把控, 拍摄角度、距离与光照条件在长达4周的干预中获得了实操性的持续微调优化。第三, AI辅助组受试者在4周训练中形成了相对稳定的动作模式与拍摄站位习惯, 减少了动作变异性对姿态估计精度的干扰。

上述现场一致性表现为低成本移动端CV方案作为VBT客观速度反馈工具的技术可行性提供了强有力的实证支撑。

\subsection{热身阶段系统性偏差的诊断与工程校正路径}
\label{sec:systematic_bias}

在Rep级现场一致性高度令人满意的同时, 热身课次级配对分析（n=88）呈现出一个具有独特方法学意义的现象：ICC依然达到0.987的优秀水平, Pearson相关系数高达0.9997, 但Bland-Altman分析揭示了+0.0212 m/s的稳定系统性正向偏移, 且88对配对中100\%均呈现App高于GA的正向偏差, 无一对呈现零偏移或负偏移。

100\%一致方向的偏移模式远远超出了随机测量误差所能解释的范畴, 指向App与GA之间存在稳定的算法口径差异, 而非测量精度问题。

关于该系统性偏差的具体机制, 研究团队在数据层面提出了两种可检验的假设。假设一为速度取值口径差异：SportSci Pro在热身阶段可能取组内所有有效重复的最大速度值作为热身速度输出, 而GymAware可能采用组内重复速度的算术均值。假设二为信号滤波策略差异：MediaPipe基于关键点位移的差分计算与GymAware基于光学编码器的直接位移测量在信号平滑窗口、噪声滤除阈值与速度峰值提取上可能采用不同参数。这两种机制假设需通过后续的工程层面代码审查与信号处理机制核对予以最终确认。

尽管偏移机制有待进一步验证, 其工程层面的实用校正路径在本研究数据基础上已具备可操作性。由于偏移量在88对配对中呈现出高度稳定的特征（均值+0.021 m/s, 标准差极小）, 可以通过线性校正实现应用层的偏差补偿：App校准值=App实测值$-$0.021 m/s。对于依赖App热身速度进行负荷--速度曲线（LVP）建模或速度--强度反推的应用场景, 建议在系统层面内置这一线性偏移校正模块；对于主要关注相对速度损失率（VL\%）而非绝对速度值的应用场景, 由于VL\%计算涉及组内相对比值, 系统性偏移会被自然抵消, 无需额外校正。

\subsection{生态效度证据的受试者代表性局限与外推边界}
\label{sec:ecological_limitations}

尽管本研究的Rep级（n=43）与热身级（n=88）配对分析在配对数量上具备充分的统计功效, 但两类数据的受试者覆盖存在明显差异。热身级配对数据来自AI辅助组全部11名受试者, 具备较强的受试者代表性；Rep级配对数据则仅来自3名受试者（P006, P016, P025）, 这一受试者覆盖限制构成了Rep级生态效度证据的核心边界。

这一局限的方法学表征体现在ICC置信区间的宽窄差异上。热身级ICC=0.987的95\%置信区间为(0.979, 0.991), 宽度仅0.012, CI精度高度收敛；Rep级ICC=0.991的95\%置信区间为(0.939, 0.992), 尽管点估计极高, 但下界扩展至0.939, 反映了受试者覆盖不足带来的抽样不确定性。

这一方法学观察对未来CV方案的效度验证研究具有明确的方法学启示：扩大受试者覆盖数量（而非仅扩大配对数量）应是提升生态效度证据外推性的优先方向。

综合本节讨论, SportSci Pro移动端视觉方案在真实抗阻训练场景中的应用边界可总结为如下三点：在典型深蹲工作负荷区间（0.4--0.8 m/s）内, 方案已达到Rep级现场一致性的可接受水平；在轻负荷热身场景（速度$>$0.7 m/s）中, 方案存在稳定的+0.021 m/s系统性正偏移, 需在系统层实施线性校正或改用相对指标；在人群外推性方面, Rep级证据主要反映的是"技术可行性"而非"人群代表性"。

\section{训练方案的可行性判定与推进决策依据}
\label{sec:feasibility_discussion}

本节回应目的2, 对第4章4.3与4.4呈现的CONSORT流程与推进标准对照进行深入解读, 并提供是否推进至大规模确证性RCT的决策依据。

\subsection{依从受试者内的高执行质量与预设标准达成}
\label{sec:high_compliance}

在进入正式训练并保持依从的24名PP受试者中, 方案的执行质量在多个维度上达到甚至超越预设推进标准。课次层训练完成率达98.4\%（189/192）；高依从率达100\%（24/24）；AI辅助组速度目标命中率达80.2\%$\pm$6.9\%, 全部11人均达$\geq$60\%预设标准；自动减组规则的触发逻辑与降组建议执行忠实度达100\%；4周干预期间未发生严重不良事件。

这些指标从执行层面证实：一旦受试者进入正式训练并保持依从状态, 本研究设计的4周非线性周期化训练方案与基于速度损失阈值的数智化监控算法在日常执行层面具备扎实的可操作性与安全性。

\subsection{保留率缺口的结构性诊断}
\label{sec:retention_gap}

从随机化后至主要后测完成的全流程视角审视, 方案的保留率呈现出显著缺口：24/36=66.7\%的主要后测完成率未达$\geq$80\%预设推进阈值。12名脱落者的时间分布呈现出鲜明的集中特征：7名（58.3\%）发生于随机化后至正式干预开始前的T0基线测试过渡阶段, 仅5名（41.7\%）发生于4周干预期间。

这一脱落时间分布提示, 保留率缺口的主要来源不是干预方案本身的负担（干预期脱落5/29=17.2\%, 处于抗阻训练干预研究的正常波动范围）, 而是基线导入流程的结构性问题。本研究采用的两测试日方案在时间上跨越48小时, 且要求受试者在两次测试日之间保持稳定的日程可用性, 这一时间跨度在在校大学生的日程冲突面前显得脆弱。

其次, T0期发生的1例非严重不良事件（P028腰部轻度肌肉拉伤）虽未导致其他受试者的连锁退出, 但P028腰部损伤虽为非严重不良事件, 其发生于随机分组后的基线测试阶段, 候选受试者处于知情同意签署后的高期待状态, 早期损伤事件在该心理状态下的间接影响不容完全排除。该效应难以在事后精确量化, 建议后续研究在知情同意环节增加对候选受试者潜在心理反应的预判说明。

第三, 干预期间缺乏针对脱落风险的主动随访干预机制。当前研究流程主要依赖受试者自主到访完成训练, 未建立每周主动随访、缺席后跟进、退出前挽留等标准化保留策略。

\subsection{系统可用性未达标的可行性影响}
\label{sec:sus_feasibility}

除保留率外, AI辅助组SUS评分65.00$\pm$4.61分未达70分预设推进阈值, 构成第二项未达标的可行性指标。11名AI辅助组受试者中仅P004一人（9.1\%）SUS评分达到70分。这一评分水平在SUS常模基准中处于"边缘可接受"区间（65--70分）, 既非完全不可用, 也远未达到用户友好水平。

值得注意的是, SUS评分与其他TAM接受度指标之间存在明显的分离, 这一二元现象将在5.7节结合质性访谈证据深入解读。就可行性判定而言, SUS未达标的直接含义是：即使受试者认可系统的功能价值并愿意持续使用, 当前原型系统的操作复杂性、学习曲线与流程流畅性仍构成阻碍其规模化推广的现实障碍。

\subsection{推进决策与流程优化路径}
\label{sec:go_no_go}

基于上述可行性证据的综合评估, 本研究对目的2的核心判定是：当前方案在依从受试者中已经具备高度的执行可行性与安全性, 但需要在招募随访流程与系统交互体验两个方面实施针对性重构后, 方可推进至大规模确证性RCT。具体而言：

在保留率优化方面, 建议实施四项结构性改进：其一, 压缩基线导入流程, 将当前的两测试日方案合并为一个测试日（约3至4小时的集中测试）, 减少受试者的时间投入与日程冲突风险；其二, 建立自随机化起的主动随访机制, 为每名受试者分配专属研究协调员, 通过每周简短随访主动了解出勤意向与障碍因素；其三, 为客观原因错过预定测试日的受试者建立灵活的补测时间窗口（如允许$\pm$3天偏离）, 避免非关键性的随访脱落；其四, 在现有交通补贴基础上增设阶段性里程碑激励, 通过完成4周训练、完成T1后测等节点的即时强化提升持续参与动力。

这一判定构成本研究作为预试验性研究的核心方法学产出。

\section{训练执行结构分化的机制解析}
\label{sec:execution_structure}

本节回应目的3的第一个子维度, 对第4章4.5呈现的"两组总外部负荷高度对等但AI辅助组总重复次数显著更少"这一执行结构分化进行机制解读。

\subsection{等剂量条件下的行为学结构分化}
\label{sec:dose_equivalent}

本研究最引人关注的训练执行发现是：尽管两组在总外部负荷（$p$=0.854）与深蹲主项负荷（$p$=0.988）上高度对等, AI辅助组却表现出显著更少的完成总重复次数（100.3 vs. 105.5次, $p$=0.035, $g$=$-$0.95）, 相当于4周内少完成约5.2次重复。这一看似微小的次数差异在等负荷的前提下意味着AI辅助组每次重复所承受的平均机械负荷高出约6.1\%, 呈现"高单位负荷、少重复次数"的负荷结构特征。

从行为学机制角度分析, 这一分化直接源于两种反馈路径的根本性调控逻辑差异。AI辅助组的"速度损失阈值截断"机制以实时监测的杠铃运动速度作为组次终止的客观判定标准；自我指导组的"RIR主观完成框架"则依托书面训练卡规定的组数$\times$次数目标, 以"保留2至3次"为主观终止判据。

\subsection{质性叙事的深度印证}
\label{sec:qualitative_印证}

质性访谈从微观行为层面为上述结构分化提供了主观机制层面的佐证。在经历过自动减组的AI辅助组受试者中, P010与P016的叙事极具说明性：两位受试者均描述了在触发自动减组时产生的"计划被打乱"与"训练刺激是否会被削减"的即时心理落差, 并在事后通过"肌肉酸痛感明显减轻"与"下一课次状态恢复更佳"的体感经验, 逐步接受了算法的安全保护逻辑。这一"初期疑虑--事后认可"的认知转向, 与AI辅助组在S4至S6负荷累积期Hooper指数持续攀升后、在S8减载期逐步恢复的先升后降轨迹形成了微观行为体验与客观恢复指标之间的互证。

相反, 自我指导组中严格内化RIR动作质量边界的P013, 其叙事展现了规范化自我指导的内在执行逻辑：P013将"出现躯干过度前倾、膝内扣或动作轨迹明显停顿"作为RIR归零的刚性边界。这一叙事表明, 规范化RIR自我指导的有效性高度依赖于训练者将抽象的主观储备概念准确转化为客观动作技术边界的能力。

反之, P015与P021呈现了规范化RIR自我指导在动作技术锚定不足时的失效模式。P015将"咬牙慢速硬顶、伴随躯干代偿完成的最后几次"依然计入完成次数, 其体能数据表现为"高最大力量增长（1RM +17.5 kg）但纯向心爆发力停滞（CMJ +0.6 cm）"的分离特征——这一现象提示, 当RIR概念未能锚定在动作技术边界上时, 规范化自我指导可能诱发疲劳性代偿完成, 从而选择性地损害对动作质量高度敏感的爆发力适应。

\subsection{操作成本的运动生理学解释}
\label{sec:operation_cost}

研究还发现AI辅助组的单次训练时长平均比自我指导组多约24分钟（822.7 vs. 798.8 min, $p$=0.072）, 这一边缘差异具有合理的运动生理学基础。首先, "高单位负荷、少重复次数"的执行结构意味着AI辅助组受试者在每组举重中需要更充分地募集运动单位；速度损失阈值的判断与确认过程增加了组间的过程性等待时间；第三, Readiness模块的训练前CMJ评估与算法决策确认进一步延长了每次训练的前置准备时间。

值得强调的是, 两组在内部负荷指标（sRPE 6.5 vs. 6.4, Hooper均值20.4 vs. 19.8）上高度对等, 表明尽管执行结构存在分化, 两组受试者在主观疲劳感知与日常恢复状态上并未展现出系统性差异。这一发现驳斥了一种可能的替代假设——即AI辅助组的算法截断会导致更低的内部负荷反应。

\section{主要体能与心理结局的解释与审慎评价}
\label{sec:outcomes_discussion}

本节回应目的3的第二个子维度, 对第4章4.7呈现的五项探索性ANCOVA结局进行机制解读与审慎评价。

\subsection{最大力量与纯向心爆发力未表现组间差异的机制分析}
\label{sec:1rm_no_diff}

本研究在深蹲绝对1RM（$p$=0.266）、相对1RM（$p$=0.273）及蹲跳SJ高度（$p$=0.622）上均未发现统计学显著的组间差异。这一结果表面上似乎出乎意料, 但从抗阻训练的神经肌肉适应机制审视, 具有深层的生理学合理性。

两组在总外部负荷上的高度对等意味着两组在4周干预期间所累积的总机械功输出基本相当。更关键的是, 两组受试者均被引导在相对接近力竭的区间执行训练：AI辅助组的速度损失阈值设定与自我指导组的RIR目标均指向相同的"接近力竭但不过度力竭"的内部刺激区间。Pareja-Blanco等（2017）的研究表明, 在等负荷条件下, 以较重负荷完成较少次数与以较轻负荷完成更多次数, 在诱发肌肉肥大方面具有相当的效应——本研究的1RM与SJ数据为该结论在力量适应维度提供了初步的预试验性佐证。

从量效关系的角度, 两组在最大力量指标上未展现差异, 还可能与4周干预周期的时长限制有关。当前4周干预周期的终点测量可能恰好捕捉到两组在神经适应层面趋于饱和、但肥大适应尚未充分展开的时间窗口——此时两组受试者均从各自路径获得相似的神经驱动效率提升, 从而在1RM指标上趋于收敛。

\subsection{CMJ方向性信号的运动学解释}
\label{sec:cmj_signal}

本研究最重要的探索性发现之一是AI辅助组在CMJ高度上呈现的边缘显著与中到大效应（调整后差值+2.16 cm, $p$=0.051, $g$=0.85, $g$ 95\% CI：0.02至1.69）。

CMJ与SJ的关键区别在于前者包含了离心预牵张阶段（stretch-shortening cycle, SSC）, 后者为纯向心收缩。Sánchez-Medina等（2011）的研究证实, SSC效率高度依赖于动作执行质量与神经激活时机——而过度的肌肉疲劳会损害SSC的弹性成分利用效率。

基于上述生理学框架, AI辅助组在CMJ上的优势可以归因于两种可能的协同机制。第一种机制涉及速度损失阈值的即时截断效应。当AI辅助组受试者在某组训练中出现明显的速度衰减时, 速度反馈促使其在离心阶段速度开始显著下降前即终止当组, 从而在后续组次中维持更高的SSC利用效率。第二种机制涉及最大意图速度（MIV）的保护。MIV指令要求受试者在每次举重中以最大意图速度完成向心阶段, 而速度损失阈值截断了低速低效的"垃圾重复"——这一策略可能使AI辅助组受试者在每次训练的后半段仍能保持较高的SSC激活质量。

然而, 必须强调的是, CMJ的组间差异在FDR校正后未达统计学显著（校正后$p$=0.126）, 差值95\% CI下界近似零, Hedges'$g$的95\% CI下界仅略高于零。因此, 上述机制解释应被定位为假设生成的探索性观察, 而非确证的因果机制。

\subsection{训练自我效能极强优势的理论阐释}
\label{sec:self_efficacy_theory}

本研究在训练自我效能上发现的超大效应（Hedges'$g$=2.70, $p<0.001$）是最具一致性与稳健性的量化发现, 且在质性叙事中获得了一致的深度印证。从理论框架角度, 这一优势可以整合Wulf等（2016）提出的OPTIMAL运动学习理论与Bandura（1997）的自我效能理论进行解释。

Wulf的OPTIMAL理论强调两种外部辅助手段对运动技能习得与表现具有最强的促进效应：一是增强期望效能——学习者对"何种努力将产生何种结果"的确信程度；二是增强自我效能——学习者对自身完成特定任务能力的信念。AI辅助组受试者所接收的客观速度量化反馈直接强化了其对"何种努力产生何种结果"的联结认知, 而即时达标信号则在每次训练中提供了稳定的掌握体验——后者正是Bandura所确认的自我效能最强来源。此外, 算法的安全边界赋予受试者对"训练过程安全可控"的信念, 进一步强化了其对完成训练任务的整体自我效能判断。

值得注意的是, AI辅助组与自我指导组在基线训练自我效能上高度均衡（65.45 vs. 65.69分）, 而经过4周差异化干预后, AI辅助组评分攀升至83.18分（净增+17.73分）, 自我指导组升至73.69分（净增+8.00分）——两组在前测高度均衡条件下产生了显著的组间分化。这一悬殊的心理效应差距提示：数智化监控的即时客观反馈所提供的掌握体验强度, 远高于规范化RIR自我指导的主观处方所带来的认知收益。

然而, 本研究在解读这一极强心理效应时必须保持三重审慎。第一, AI辅助组受试者知晓自身所使用的是"数智化监控系统", 开放标签设计不可避免地引入了期望效应的混杂。第二, 训练自我效能ANCOVA模型残差Shapiro-Wilk检验$p$=0.001, 偏离正态假设。第三, 训练自我效能量表本身对高分区间可能存在天花板效应, AI辅助组均值83.18分已接近量表上限, 量表刻度分辨率对超大效应的估计可能存在膨胀风险。三项审慎的应对策略见5.8节与5.9节。

\section{Hooper恢复轨迹与内部负荷响应的机制解读}
\label{sec:hooper_discussion}

本节回应目的3的第三个子维度, 对第4章4.6呈现的Hooper LMM分析结果进行机制解读, 并阐释组别与课次两种效应格局的方法学与实践意义。

\subsection{先升后降轨迹的训练学理机制解读}
\label{sec:up_down_trajectory}

本研究通过线性混合效应模型对8次训练课前的Hooper主观恢复总分进行了纵向分析, 取得了三项相互印证的关键发现：组别主效应不显著（$F$(1, 22.0)=0.67, $p$=0.421）、组别与课次交互效应不显著（$F$(7, 154.0)=0.32, $p$=0.945）, 但课次主效应达到高度统计学显著（$F$(7, 154.0)=3.09, $p$=0.004）。

这一先升后降的时间演变形态与本研究采用的4周训练周期化设计（第1周导入适应$\to$第2至3周超负荷$\to$第4周减载恢复）呈现出高度精准的对应关系。在S4至S6峰值阶段, 训练主项组数从第1周的4组增加至第2至3周的5组, 训练日A的相对负荷从70\%--75\% 1RM递增至80\%--85\% 1RM——这一累积性负荷升级直接反映在Hooper总分攀升上。在S8减载阶段, 主项组数缩减至3组, 主观恢复状态显著改善并出现超量恢复效应——这与经典抗阻训练周期化理论中"减载周期通过降低急性疲劳累积促进适应性超量恢复"的机制预期完全吻合。

由于本研究未对各课次之间的两两pairwise比较进行显著性检验, 上述时间演变形态属于对课次主效应显著（$p$=0.004）的形态学描述, 不构成对特定课次间差异的统计学推断。

\subsection{两组恢复轨迹平行与"等应激假设"的证实}
\label{sec:equal_stress}

在课次主效应显著的同时, 组别与课次的交互效应高度不显著（$F$(7, 154.0)=0.32, $p$=0.945）, 跨全部8个课次的组间Hooper差值均值仅为0.62分（95\% CI：$-$0.95至2.20）, 全部8个课次的组间差值95\%置信区间均跨越无效值0。

从方法学角度审视, 两组恢复轨迹的高度平行性有力驳斥了两种可能的替代假设。替代假设一认为, AI辅助组的自动减组与速度损失阈值截断可能因"训练量削减"而导致更低的疲劳应激反应；替代假设二认为, AI辅助组因每次重复的机械负荷更高（+6.1\%）可能诱发更强的疲劳应激。本研究的数据表明, 无论从减少刺激还是加重刺激的假设出发, 都无法在4周干预期间的Hooper主观恢复轨迹上观察到组间分化。这一"等应激现象"从内部负荷反应维度补充证实了两种反馈路径在等外部剂量下均能达到相似的内部刺激区间, 为后续两组适应效应对比的方法学公平性提供了关键的过程性证据。

从实践应用角度, 两组恢复轨迹平行也为训练者与教练员提供了一个安心的实证保证：引入数智化速度反馈与算法安全边界并未在训练者身上诱发额外的主观应激。

\subsection{Kenward-Roger自由度修正的方法学价值}
\label{sec:kr_methodology}

本研究在Hooper LMM分析中采用了Kenward-Roger（KR）Type III F检验作为固定效应显著性推断的核心方法, 这一选择具有超出本研究单一发现的普遍方法学启示价值。

KR修正的必要性源于小样本纵向重复测量设计的固有统计学挑战。在本研究这样每组仅11至13名受试者、每人8次重复测量的小样本纵向设计中, 渐近Wald卡方检验会系统性地低估残差方差的不确定性并高估自由度, 从而导致统计功效损失。

在本研究的具体情境中, 若采用传统Wald卡方检验, Hooper课次主效应的$p$值约为0.119, 将被判定为不显著；而采用KR F检验后, $p$值降至0.004, 充分捕获了先升后降时间轨迹这一具有明确训练学意义的效应。

这一方法学发现对未来运动科学领域小样本纵向研究的统计推断具有重要启示。本研究建议后续开展类似小样本纵向RCT的研究者, 在使用R语言的\texttt{lmerTest::lmer()}时优先选择\texttt{ddf="Kenward-Roger"}参数, 或在其他统计平台主动指定Satterthwaite或Kenward-Roger自由度修正, 以避免因方法学工具选择而产生的效应遗漏。

\section{系统可用性、人机信任与用户自主权的动态平衡}
\label{sec:acceptance_discussion}

本节回应目的4, 对第4章4.9与4.10呈现的用户接受度证据进行整合性机制解读。

\subsection{"高价值认同vs.易用性瓶颈"二元矛盾的深层解析}
\label{sec:usability_paradox}

本研究在系统评价维度取得的最引人深思的发现是AI辅助组量化评价数据中呈现的鲜明二元分离特征：一方面, 感知有用性（PU=4.41$\pm$0.20）、信任度（Trust=4.41$\pm$0.20）与持续使用意愿（Intention=4.49$\pm$0.27）三项TAM核心接受度维度的均分均超过4.4分（满分5分, 得分率超过88\%）, 呈现出令人瞩目的高价值认同；另一方面, 系统可用性量表（SUS）得分均值仅为65.00$\pm$4.61分, 11名受试者中仅P004一人（9.1\%）达到70分推进阈值。

这一表面矛盾的量化数据格局在质性访谈中得到了深度机制层面的破解。系统的功能价值获得压倒性认同——受试者P004、P025高度评价实时速度显示"消除了以往加重前的犹豫", P010、P016认为算法自动减组"提供了明确的安全保护", P018承认"数据本身很有价值"。然而, 系统的过程性操作体验则暴露了显著的摩擦点：设备连接的等待时间、组间的实时数据确认、专业术语的认知负荷（MCV、VL\%、$\Delta$v等VBT专业指标）、训练节奏的中断。这些过程性摩擦解释了SUS评分为何仅为65.00分。

这一二元矛盾在原型系统评估领域并不罕见, 其本质是"核心价值验证"与"交互体验成熟度"两个开发阶段的时间错位。本研究的SportSci Pro已在核心速度测量、算法决策与安全保护等功能层面达到了可接受的成熟度, 但在交互设计层面尚处于早期原型阶段。

\subsection{用户自主权与算法控制的动态博弈：分层匹配框架的实证支撑}
\label{sec:autonomy_algorithm}

质性访谈发现了一个远超简单"接受与拒绝"二分法的深层现象：用户对数字化系统的最优互动模式受个体训练经验、力量水平与自主性倾向的显著调节。

从自我决定理论（Self-Determination Theory, SDT；Ryan \& Deci, 2017）视角审视, 人类行为动机的持续性依赖于三种基本心理需要的满足：胜任感、自主感与归属感。数字化训练系统对这三种心理需要的影响并非单一方向。

"积极接受型"个案（P004、P025）代表了系统SDT需要满足的最优场景。这类受试者训练经验相对有限, 客观速度反馈提供的即时达标信号直接增强了胜任感；算法明确的负荷建议将有限的自主选择空间转化为对训练本身的深度专注。

"安全依赖型"个案（P010、P016）呈现出算法控制在特殊人群中的独特价值。Readiness自动减组机制在其主观感知层面转化为外部保护性权威——将"是否应该继续训练"这一心理压力从个人决策转移至客观算法。

"批判接受型"个案（P018、P012、P014）则揭示了算法控制与用户自主权之间的核心张力点。这类受试者具备扎实的训练经验与稳定的本体感觉基础, 当算法的自动化决策不再是"能力补充", 反而在客观上削弱了他们的自主感。

这一分层匹配现象引出了对未来数字化训练工具设计的核心方法学建议：面向大众市场的数智化训练系统不能采用"一刀切"的自动化模式, 而应设计"分级自主授权"的自适应交互框架。

\subsection{技术反馈依赖风险的审慎警示}
\label{sec:feedback_dependency}

在讨论数字化训练系统的价值时, 本研究认为必须客观承认技术介入可能带来的潜在负面效应, 避免"技术乐观主义"的过度延伸。

Winstein与Schmidt（1990）提出的引导假说（guidance hypothesis）指出, 过于频繁与强烈的外部反馈虽然可以在训练当下改善动作表现, 但长期可能形成"反馈依赖"——学习者变得只能在外部反馈存在时表现良好, 一旦反馈缺失即出现表现衰退。

Jakicic等（2016）在《美国医学会杂志》（JAMA）发表的24个月随机对照试验为技术反馈依赖风险提供了直接的临床证据：与对照组相比, 添加可穿戴健身追踪设备的干预组在24个月后的减重效果反而更差（$-$3.5 kg vs. $-$5.9 kg, $p$=0.04）。

尽管本研究的干预周期（4周）远短于Jakicic研究（24个月）, Jakicic研究提示的长期反馈依赖风险仍值得本研究严肃对待。本研究提示后续正式RCT与商业化系统部署应考虑以下三项反馈依赖风险的缓解策略：其一, 在长期使用中引入间歇性反馈撤除机制（如每月安排1周的"无反馈自主训练周"）；其二, 在系统设计中显性区分"辅助模式"与"独立模式"；其三, 将RIR判断与本体感觉训练内嵌于数字化系统的教学模块中。

\section{研究局限性}
\label{sec:limitations}

本研究在设计上虽严格遵循CONSORT-Pilot声明与CERT报告规范, 但作为预试验性混合方法研究, 其固有的小样本特征与开放标签设计仍构成了多项系统性局限, 需要在解读本研究结论时进行审慎的边界标注。

\subsection{样本规模与主要结局的统计学功效约束}
\label{sec:sample_power}

本研究PP样本量为24名（AI辅助组11名, 自我指导组13名）, 这一样本规模在预试验性研究中具备基本的过程可行性检验与假设生成功能, 但在统计学功效层面对组间中等效应（$d$=0.50）的检测功效仅约为0.33——远低于正式RCT通常要求的0.80功效水平。这意味着本研究在主要体能结局指标上存在较高的II类错误风险。

以CMJ高度这一呈现边缘信号的结局为例, 本研究检测到的中到大效应（$g$=0.85）在当前样本下仅达到边缘显著（$p$=0.051）, FDR校正后失去统计显著性（校正后$p$=0.126）。若该效应真实存在, 在扩展至每组24人（共48人）的正式RCT样本后, 该效应在$\alpha$=0.05双侧检验下可达到约0.80的统计学功效。

\subsection{干预周期对力量适应时程的限制}
\label{sec:intervention_duration}

本研究的总干预周期仅为4周, 这一设计时长在训练学研究的方法学框架中处于短期预试验的范畴。根据抗阻训练适应时程的阶段性规律, 神经肌肉系统对外部负荷刺激的适应遵循可辨识的时序模式：神经驱动效率的早期改善通常在2至4周内可观测到, 而肌肉横截面积增加与肌纤维肥大成分的实质性贡献通常需要8至12周以上方可达到具有实际意义的量级。

本研究4周的终点测量恰好位于"神经适应趋于饱和、肥大适应尚未充分展开"的时间窗口, 两组受试者无论从何种反馈路径出发, 均在神经驱动的早期适应阶段获得了相似的收益——这一时程特征可能是1RM与SJ指标未呈现组间差异的另一重要原因。

\subsection{开放标签设计与期望效应混杂}
\label{sec:blinding}

本研究采用开放标签设计, 两组受试者均知晓自身所属的干预条件, 研究人员也全程知晓各受试者的组别归属。这一设计选择对于评估"数智化监控系统在真实使用场景中的表现"具有内在的必要性, 但不可避免地引入了期望效应的系统性混杂。

在训练自我效能维度, 这一混杂尤为关键。AI辅助组受试者因"使用了高科技智能系统"而产生的心理期待与自我暗示, 可能部分解释了为何AI辅助组的自我效能增量远超自我指导组。换言之, 自我效能的极强组间优势可能混合了速度反馈的技术性效应与"高科技组"的安慰剂效应, 两者无法在当前设计框架下有效分离。后续正式RCT应通过设立注意力等量对照组来控制期望效应。

\subsection{生态效度证据的受试者代表性局限}
\label{sec:ecological_lim}

尽管本研究热身级配对分析已覆盖AI辅助组全部11名受试者（n=88）, 代表性较强, 但Rep级配对分析（n=43）仅来自3名受试者（P006, P016, P025）——这一Rep级证据的受试者覆盖限制意味着当前Rep级ICC、Bias与MAE参数估计实质上仅代表了这3名受试者在特定拍摄环境、动作模式与训练习惯下的系统表现, 尚无法支撑"任意AI辅助组受试者在任意场景下使用App均可获得同等Rep级一致性"的外推结论。

后续正式RCT应将App与GA的同步记录纳入每名AI辅助组受试者的标准Rep级采样流程, 并将受试者的身体特征多样性与训练环境多样性纳入效度证据的代表性格局设计。

\subsection{质性分析的单人编码局限}
\label{sec:qualitative_lim}

本研究的质性编码工作由单一研究者完成, 未采用独立双人编码与编码者间一致性检验的标准化质性信度保障程序。此外, 本研究的质性抽样存在一个明显的覆盖局限：实际纳入访谈的AI辅助组5名受试者均来自低力量层, 高力量层未覆盖；因此本研究关于"高力量层AI使用体验"的质性证据完全依赖自我指导组的P012与P014做间接推测。后续开展更大规模混合方法研究时, 建议引入至少一名独立编码者进行双盲编码, 并通过Kappa系数等指标量化编码者间一致性, 同时确保AI辅助组高力量层受试者的完整覆盖。

\subsection{训练自我效能量表的天花板效应风险}
\label{sec:se_ceiling}

如5.5.3节所述, 本研究训练自我效能量表的组间超大效应（$g$=2.70）需在天花板效应的框架下审慎解读。AI辅助组T1均值达83.18分, 若量表理论上限接近90至100分区间, 则该组的进一步增长空间已相对有限；此外, ANCOVA模型残差Shapiro-Wilk $p$=0.001, 明显偏离正态假设, 提示当前效应量估计可能受量表刻度分辨率与残差分布形态的双重影响。后续正式RCT建议在训练自我效能测量上引入更高刻度分辨率的量表版本（如从0至100的连续视觉模拟评分）, 并采用秩变换ANCOVA或bootstrap校验的稳健分析方法。

\section{后续正式随机对照试验的设计建议与实践启示}
\label{sec:rct_recommendations}

基于本研究在流程可行性、测量效度、训练执行、适应信号、系统接受度与质性机制方面所积累的全部实证经验, 本节提出面向后续正式随机对照试验的设计参数建议、系统工程迭代方向以及面向大众训练实践的审慎建议。

\subsection{正式RCT的样本量估算与参数依据}
\label{sec:sample_size}

后续正式RCT的样本量估算应以CMJ高度作为主要结局指标。以CMJ调整后组间差值+2.16 cm为预期效应量参数、PP样本CMJ合并标准差约2.54 cm为噪声参数, 双侧$\alpha$=0.05、目标统计学功效1-$\beta$=0.80进行两独立组比较的样本量估算, 所需每组最小样本量约为20人；考虑约20\%的随访脱落率, 每组需随机化约24人, 总计48人。若加入注意力等量对照组以控制开放标签设计带来的期望效应混杂, 则总样本量需扩展至约72人。建议在期中分析时根据累积数据进行样本量重新估算, 以根据实际效应量调整总样本目标。上述估算基于本研究PP样本的效应量估计, 精确数字需在后续预试验或文献荟萃中更新。

\subsection{注意力等量对照组的设立}
\label{sec:attention_control}

鉴于开放标签设计中自我效能结局存在期望效应混杂的内在风险, 后续正式RCT建议在传统"AI辅助组"与"自我指导组"二分设计的基础上, 增设第三臂——注意力等量对照组（Attention Control Group, AC组）。AC组受试者接受等量的额外关注与流程参与（每周1次与研究团队的电话/短信随访、每次训练后填写简短的训练感受日志）, 但其训练指导内容不包含实时速度反馈, 仅提供与自我指导组相同的规范化RIR书面训练卡。

AC组的引入可使研究者能够通过组间比较将"数智化速度反馈的技术效应"从"额外关注与流程参与的心理效应"中统计分离。若三臂设计在招募可行性上存在困难, 也可采用2$\times$2因子设计（实时反馈：有/无$\times$关注强度：高/低）, 以更高效地解析两种效应的交互作用。

\subsection{招募保留率提升的多层次策略}
\label{sec:retention_strategy}

本研究66.7\%的保留率揭示了招募随访流程中的系统性薄弱环节。建议从以下四个维度系统性提升保留率。第一, 压缩基线导入流程：将分散在多日的基线信息采集压缩至单日集中完成, 减少首次接触至随机化之间的等待时间窗口。第二, 实施主动随访机制：自随机化之日起, 每名受试者分配专属研究协调员, 建立每周1次的简短随访联系, 主动了解受试者的出勤意向与障碍因素。第三, 建立灵活的补测安排：对于因客观原因错过预定测试日的受试者, 在统计方法支持范围内提供灵活的补测时间窗口（如允许$\pm$3天的偏离）。第四, 强化激励结构：在现有交通补贴基础上增设里程碑激励（如完成4周全部训练课次后的一次性奖励）, 通过阶段性目标的即时强化提升受试者的持续参与动力。

\subsection{系统交互体验的工程迭代方向}
\label{sec:system_iteration}

基于SUS量化发现与质性访谈所报告的操作摩擦, 本研究提出以下明确的系统交互优化优先级清单。第一优先级：极简UI与数据可视化重构。将现有界面从"多指标并列展示"重构为"单核心数字+可选详情面板"的极简模式, 以颜色区间条（绿色=达标、黄色=接近阈值、红色=已衰减）与直观图形替代MCV、VL\%、$\Delta$v等数值, 降低非专业用户的认知门槛。第二优先级：自动化设备连接。将设备就绪时间从当前2至3分钟/次压缩至$<$30秒。第三优先级：算法反馈的可解释性增强。将当前"减组通知"的机器通知式输出改为自然语言解释式输出, 同步告知用户"为什么"以缓解算法干预引发的刺激焦虑。第四优先级：热身阶段系统性偏移的工程层校正模块。在系统架构中内置App校准值=App实测值$-$0.021 m/s的线性补偿, 确保在进行LVP建模时输出值的绝对准确性。

\subsection{大众训练实践的分层应用建议}
\label{sec:practice_recommendations}

基于5类行为特征模式与SDT分层匹配框架的实证证据, 本研究对数智化监控工具向大众训练场景的下沉提出分层应用的审慎建议, 不倡导"一刀切"的全人群推广。

对于训练经验不足1年的初学者, P015与P021的质性案例表明, 初学者难以准确将RIR概念转化为动作技术边界, 极易在追求大重量过程中累积疲劳与技术变形。对于这一群体, 引入基于速度损失阈值的客观反馈具有明确的应用价值, 可作为过渡性辅助工具, 推荐以"入门级模式"部署。

对于训练年限在1至3年的进阶训练者, 本研究提供了两种反馈路径均可实现有效适应的证据。对于这一群体, 反馈工具的价值不在于教导基础动作, 而在于提供数据驱动的训练监控与疲劳诊断, 推荐以"进阶级模式"部署, 强调自主确认与数据参考。

对于训练年限超过3年的高水平训练者（以P012、P014为典型代表）, 数智化工具应定位为"高精度数据仪表盘"而非"训练指挥官"。需要说明的是, 该层级的行为特征证据主要来自本研究自我指导组的高经验受试者, 其在AI辅助条件下的实际接受模式仍待后续研究直接验证。

上述三层分级框架与5.7.2节归纳的5类行为特征模式相互印证, 建议后续正式RCT在系统功能设计时即引入上述分层授权逻辑, 同步评估不同分层模式对接受度和训练适应的影响。
